%%%%%%%% ICML 2026 EXAMPLE LATEX SUBMISSION FILE %%%%%%%%%%%%%%%%%

\documentclass{article}

% Recommended, but optional, packages for figures and better typesetting:
\usepackage{microtype}
\usepackage{graphicx}
\usepackage{subcaption}
\usepackage{booktabs} % for professional tables

% hyperref makes hyperlinks in the resulting PDF.
% If your build breaks (sometimes temporarily if a hyperlink spans a page)
% please comment out the following usepackage line and replace
% \usepackage{icml2026} with \usepackage[nohyperref]{icml2026} above.
\usepackage{hyperref}


% Attempt to make hyperref and algorithmic work together better:
\newcommand{\theHalgorithm}{\arabic{algorithm}}

% Use the following line for the initial blind version submitted for review:
\usepackage{icml2026}

% For preprint, use
% \usepackage[preprint]{icml2026}

% If accepted, instead use the following line for the camera-ready submission:
% \usepackage[accepted]{icml2026}

\usepackage{amsmath}
\usepackage{amssymb}
\usepackage{mathtools}
\usepackage{amsthm}


% if you use cleveref..
\usepackage[capitalize,noabbrev]{cleveref}

%%%%%%%%%%%%%%%%%%%%%%%%%%%%%%%%
% THEOREMS
%%%%%%%%%%%%%%%%%%%%%%%%%%%%%%%%
\theoremstyle{plain}
\newtheorem{theorem}{Theorem}[section]
\newtheorem{proposition}[theorem]{Proposition}
\newtheorem{lemma}[theorem]{Lemma}
\newtheorem{corollary}[theorem]{Corollary}
\theoremstyle{definition}
\newtheorem{definition}[theorem]{Definition}
\newtheorem{assumption}[theorem]{Assumption}
\theoremstyle{remark}
\newtheorem{remark}[theorem]{Remark}

% Todonotes is useful during development; simply uncomment the next line
%    and comment out the line below the next line to turn off comments
%\usepackage[disable,textsize=tiny]{todonotes}
\usepackage[textsize=tiny]{todonotes}

% The \icmltitle you define below is probably too long as a header.
% Therefore, a short form for the running title is supplied here:
\icmltitlerunning{UEWC-Merge: Teacher-Free Test-Time Model Merging via Unsupervised EWC}

\begin{document}

\twocolumn[
  \icmltitle{Unsupervised Elastic Weight Consolidation for \\
    Teacher-Free Test-Time Model Merging}

  % It is OKAY to include author information, even for blind submissions: the
  % style file will automatically remove it for you unless you've provided
  % the [accepted] option to the icml2026 package.

  % List of affiliations: The first argument should be a (short) identifier you
  % will use later to specify author affiliations Academic affiliations
  % should list Department, University, City, Region, Country Industry
  % affiliations should list Company, City, Region, Country

  % You can specify symbols, otherwise they are numbered in order. Ideally, you
  % should not use this facility. Affiliations will be numbered in order of
  % appearance and this is the preferred way.
  \icmlsetsymbol{equal}{*}

  \begin{icmlauthorlist}
    \icmlauthor{Firstname1 Lastname1}{equal,yyy}
    \icmlauthor{Firstname2 Lastname2}{equal,yyy,comp}
    \icmlauthor{Firstname3 Lastname3}{comp}
    \icmlauthor{Firstname4 Lastname4}{sch}
    \icmlauthor{Firstname5 Lastname5}{yyy}
    \icmlauthor{Firstname6 Lastname6}{sch,yyy,comp}
    \icmlauthor{Firstname7 Lastname7}{comp}
    %\icmlauthor{}{sch}
    \icmlauthor{Firstname8 Lastname8}{sch}
    \icmlauthor{Firstname8 Lastname8}{yyy,comp}
    %\icmlauthor{}{sch}
    %\icmlauthor{}{sch}
  \end{icmlauthorlist}

  \icmlaffiliation{yyy}{Department of XXX, University of YYY, Location, Country}
  \icmlaffiliation{comp}{Company Name, Location, Country}
  \icmlaffiliation{sch}{School of ZZZ, Institute of WWW, Location, Country}

  \icmlcorrespondingauthor{Firstname1 Lastname1}{first1.last1@xxx.edu}
  \icmlcorrespondingauthor{Firstname2 Lastname2}{first2.last2@www.uk}

  % You may provide any keywords that you find helpful for describing your
  % paper; these are used to populate the "keywords" metadata in the PDF but
  % will not be shown in the document
  \icmlkeywords{Machine Learning, ICML}

  \vskip 0.3in
]

% this must go after the closing bracket ] following \twocolumn[ ...

% This command actually creates the footnote in the first column listing the
% affiliations and the copyright notice. The command takes one argument, which
% is text to display at the start of the footnote. The \icmlEqualContribution
% command is standard text for equal contribution. Remove it (just {}) if you
% do not need this facility.

% Use ONE of the following lines. DO NOT remove the command.
% If you have no special notice, KEEP empty braces:
\printAffiliationsAndNotice{}  % no special notice (required even if empty)
% Or, if applicable, use the standard equal contribution text:
% \printAffiliationsAndNotice{\icmlEqualContribution}

\begin{abstract}
  Model merging has emerged as a cost-effective paradigm to integrate specialized capabilities of independently fine-tuned expert models into a single multi-task system without requiring prohibitive joint retraining. To adapt to out-of-distribution shifts and non-stationary test streams, test-time adaptation (TTA) of model merging has recently been proposed. However, existing TTA model merging methods suffer from a severe trade-off: they are either \textit{teacher-guided} (e.g., SATA, EWC-TTA), requiring large expert models to be hosted in memory as teachers at inference time, which defeats the primary memory-saving motivation of model merging; or they are \textit{teacher-free} but highly restricted (e.g., S2C-Merge), freezing classification heads entirely to prevent decision boundary collapse, which severely limits their capability to adapt to domain corruptions. In this paper, we propose \textbf{UEWC-Merge} (Unsupervised Elastic Weight Consolidation for Teacher-Free Test-Time Model Merging), a framework that resolves this fundamental trade-off. UEWC-Merge pre-computes the diagonal Fisher Information Matrix (FIM) of each expert classification head offline, and applies an unsupervised EWC penalty during TTA to anchor the heads. This allows for the simultaneous, safe adaptation of both the layer-wise model merging coefficients and the classification heads under purely self-supervised objectives (prediction entropy minimization and task-aware consistency regularization), with zero teacher model memory overhead. Our comprehensive evaluation on a multi-task vision benchmark demonstrates that UEWC-Merge completely prevents decision boundary collapse, significantly outperforming teacher-free baselines that freeze classification heads, and matching or exceeding heavy teacher-guided methods.
\end{abstract}

\section{Introduction}
Deep neural networks have achieved remarkable success across a wide range of machine learning tasks, largely driven by pre-training on massive datasets and subsequently fine-tuning on specialized downstream domains. This workflow frequently results in highly specialized, task-specific expert checkpoints. While these expert models achieve state-of-the-art performance on their respective target domains, hosting each expert model independently at inference time incurs prohibitive VRAM, storage, and hosting overhead. This parameter bottleneck becomes especially severe when deploying models on resource-constrained commodity edge hardware, such as smart sensors, mobile devices, autonomous vehicles, or decentralized robotics. Consequently, model merging has emerged as a highly popular, computationally efficient, and cost-effective paradigm to integrate the specialized capabilities of independently fine-tuned expert models into a single, unified multi-task system without requiring prohibitive and expensive joint retraining \cite{Wortsman22,Matena22}.

Traditional static merging methods, such as Task Arithmetic \cite{Ilharco22} and TIES-Merging \cite{Yadav23}, address this by applying fixed linear or coordinate-wise operations to combine expert parameters. While highly efficient, static merging frequently suffers from severe parameter interference and representation collapse when experts have incompatible geometries, or when the merged model is exposed to shifted, out-of-distribution (OOD) test environments. To address these limitations, adaptive test-time model merging (TTMM) frameworks have recently been proposed \cite{Yang24,Jung25}. These frameworks dynamically adjust merging weights on unlabeled test-time data streams on-the-fly, allowing the model to adapt to local distribution shifts, environmental corruptions, and non-stationary streams. Recent state-of-the-art methods, such as SATA \cite{Sata26} and EWC-TTA \cite{Ewc26}, optimize layer-wise merging coefficients and task-specific classification heads dynamically on the target stream.

However, existing adaptive test-time model merging methods are fundamentally restricted by what we term the \textbf{Teacher-Overhead Paradox}: to guide online adaptation without access to ground-truth labels, they rely on \textit{teacher-guided} self-labeling or distillation objectives computed against the original, unmerged expert models. Consequently, these methods require loading and maintaining all expert models in memory as teachers at inference time. This quadruples or quintuples the VRAM and hosting footprint, directly defeating the primary resource-saving motivation of model merging and rendering TTA unfeasible on commodity edge hardware.

To break this paradox, teacher-free test-time model merging frameworks, such as S2C-Merge \cite{S2C26}, have been introduced. S2C-Merge utilizes purely self-supervised objectives, including prediction entropy minimization and consistency regularization, to adapt merging coefficients on-the-fly with zero teacher model memory overhead. However, S2C-Merge and related teacher-free methods suffer from two critical, unaddressed bottlenecks: (1) \textbf{Non-Stationary Parameter Interference:} On non-stationary, task-alternating streams, updating a single global set of merging coefficients sequentially across tasks causes severe gradient interference. Updates computed for one task overwrite and degrade parameters learned for other tasks, resulting in high parameter oscillations and negative transfer. (2) \textbf{Decision Boundary Collapse:} Optimizing both the high-dimensional feature representation and task classifiers under purely unsupervised, self-supervised objectives leads to severe decision boundary collapse. S2C-Merge resolves this only by freezing the classification heads entirely. This severely restricts their capability to adapt to domain corruptions (e.g., severe noise, blur, or contrast changes) that alter the representation distribution.

In this work, we propose \textbf{UEWC-Merge} (Unsupervised Elastic Weight Consolidation for Teacher-Free Test-Time Model Merging), a framework that completely resolves these fundamental trade-offs. To systematically address the challenges of adaptive test-time model merging under non-stationary and corrupted streams, we formulate and address three core research questions:
\textbf{RQ1 (Task-Adaptive Routing Effectiveness):} Can an ultra-lightweight task classifier achieve high robustness and guide accurate routing under severe out-of-distribution shifts without relying on heavy teacher models?
\textbf{RQ2 (Multi-Task Adaptation and Collapse Prevention):} Can our unsupervised EWC penalty effectively stabilize joint adaptation of classification heads and merging coefficients, preventing decision-boundary collapse under self-supervised objectives?
\textbf{RQ3 (Generalization and Resource Efficiency):} Does UEWC-Merge outperform existing teacher-free and teacher-guided baselines under diverse domain corruptions while achieving a significant reduction in memory and computational overhead?

Our key contributions are summarized as follows:
\begin{itemize}
    \itemsep0em
    \item We introduce \textbf{Task-Adaptive Test-Time Model Merging (TA-TTMM)}, a paradigm that completely resolves non-stationary parameter interference. We pre-train a robust, ultra-lightweight task classifier (1.3K parameters, ~5 KB) offline using robust multi-style data augmentations, enabling on-the-fly task routing of unlabeled batches with $>82\%$ accuracy even under severe corruptions.
    \item We maintain task-specific layer-wise merging coefficients $\Lambda_k$ routed on-the-fly and optimized independently, completely isolating tasks and preventing cross-task gradient interference.
    \item We pre-compute clean diagonal Fisher Information Matrix (FIM) priors for each expert's classification head offline using a tiny clean validation set (e.g., $N=200$ samples), requiring negligible storage overhead and zero teacher memory at test time.
    \item We provide a rigorous Bayesian formulation showing that our unsupervised EWC penalty acts as a Laplace approximation of the clean posterior distribution, anchoring classification heads during self-supervised TTA. This prevents decision boundary collapse while permitting classifiers to adapt and align with corrupted, out-of-distribution features.
    \item Through rigorous empirical evaluation on a multi-task vision benchmark under diverse corruptions, we demonstrate that UEWC-Merge completely stabilizes head adaptation, significantly outperforming teacher-free baselines that freeze classification heads, and decisively outperforming heavy teacher-guided baselines while saving 4$\times$ memory.
\end{itemize}

\section{Related Work}
\subsection{Model Merging and Fusion}
Weight averaging and deep model fusion combine independent neural networks sharing a common pre-trained initialization to achieve multi-task capabilities without retraining \cite{Wortsman22,Matena22}. Task Arithmetic \cite{Ilharco22} adds task vectors, while TIES-Merging \cite{Yadav23} resolves sign disagreements and parameter redundancies. RegMean \cite{jin2022regmean} optimizes linear merging of weights by minimizing L2 distance between representations, and other coordinate-wise operations have been proposed to combine expert parameters. Git-Re-Basin \cite{ainsworth2022git} addresses permutation symmetries in weight spaces to align models before merging. Dynamic model merging adapts coefficients during test-time to specialize the network to incoming data \cite{Yang24,Jung25}. However, static merging cannot handle non-stationary target streams or local domain corruptions, motivating our dynamic, adaptive paradigm.

\subsection{Model Merging versus Multi-Task Learning}
An alternative to model merging is Joint Multi-Task Learning (MTL) or Mixture-of-Experts (MoE) training, where a single network is trained simultaneously on all target datasets. While highly effective, joint MTL is computationally expensive and introduces severe logistical and privacy constraints, as it requires simultaneous access to all clean datasets during training. Furthermore, extending a trained multi-task model to new tasks requires expensive joint retraining to avoid catastrophic forgetting. In contrast, model merging represents a highly modular, post-hoc alignment paradigm. Specialized expert models can be trained independently, in decentralized environments, and merged zero-shot at deployment time. Our test-time adaptive model merging framework extends this modularity to dynamic environments, allowing the merged model to adapt to target streams post-hoc with zero joint training overhead.

\subsection{Test-Time Adaptation (TTA)}
TTA adjusts a pre-trained model to target domains using unlabeled streams on-the-fly \cite{Wang21,Liang20}. Classical TTA methods (e.g., Tent \cite{wang2021tent}) minimize prediction entropy, while consistency-based methods (e.g., CoTTA \cite{wang2022continual}) preserve knowledge under continuous shifts. S2C-Merge \cite{S2C26} combines entropy minimization and task-aware consistency to adapt merging weights. SATA \cite{Sata26} uses Fisher-guided SAM, while EWC-TTA \cite{Ewc26} precomputes FIM to regularize heads. However, both SATA and EWC-TTA require expert models as online teachers, suffering from heavy memory footprints and massive latency.

\subsection{Continual Learning and Elastic Weight Consolidation}
Continual Learning (CL) addresses catastrophic forgetting in sequential training \cite{delange2021continual}. Regularization-based CL methods, such as Elastic Weight Consolidation (EWC) \cite{kirkpatrick2017overcoming} and Synaptic Intelligence \cite{zenke2017continual}, mitigate forgetting by constraining parameters using the diagonal Fisher Information Matrix (FIM) computed on old tasks. In this work, we creatively repurpose EWC from sequential training to self-supervised TTA, utilizing it as an unsupervised decision-boundary anchor that prevents classifier collapse while permitting localized adaptation under severe domain shifts.

\section{Methodology}
Let $\Theta_{\text{pre}}$ denote the parameters of a shared pre-trained base encoder. We consider $K$ distinct tasks, producing $K$ expert encoders $\Theta_1, \dots, \Theta_K$ and their respective task classification heads $\phi_1, \dots, \phi_K$. The task vector for expert $k$ is defined as $\underline{\tau}_k = \Theta_k - \Theta_{\text{pre}}$.

\subsection{Robust Task Routing via Tiny Classifier}
To resolve task conflict on non-stationary streams without accessing ground-truth labels or loading heavy teacher models, we introduce a robust, ultra-lightweight task classifier $\psi$. The classifier $\psi$ is a single-layer CNN (1.3K parameters) trained offline on a small validation set with random Gaussian noise, Gaussian blur, and contrast corruptions. At test-time, for an incoming unlabeled batch $X_t$, we predict the task ID $\hat{k}$ via majoritarian voting across the batch:
\begin{equation}
    \hat{k} = \operatorname{argmax}_{c} \sum_{x \in X_t} \mathbb{I}[\psi(x) = c]
\end{equation}
where $\mathbb{I}[\cdot]$ represents the indicator function. The predicted task $\hat{k}$ is used on-the-fly to route $X_t$ to its task-specific merging coefficients and classification head. This batch-level routing acts as a powerful denoising filter, as voting filters out isolated sample-level classification errors under heavy noise.

\subsection{Task-Adaptive Layer-wise Model Merging}
We define task-specific, layer-wise merging coefficients \, $\Lambda_k \in \mathbb{R}^{L \times K}$ for each task $k$. The merged model's weights for parameter tensor $l$ under predicted task $\hat{k}$ are reconstructed dynamically on-the-fly using a clamp constraint:
\begin{equation}
    \Theta_{\text{merged}}^{(l)} = \Theta_{\text{pre}}^{(l)} + \sum_{m=1}^K \lambda_{\hat{k}, m}^{(l)} \underline{\tau}_m^{(l)}
\end{equation}
where $\lambda_{\hat{k}, m}^{(l)} = \operatorname{clamp}(\tilde{\lambda}_{\hat{k}, m}^{(l)} + 1/K, 0, 1)$ represent the active merging coefficients reconstructed from the raw learnable parameters $\tilde{\lambda}_{\hat{k}, m}^{(l)}$. This clamp parameterization provides stable gradients and restricts coefficients to a mathematically valid range $[0, 1]$, outperforming Softmax-based constraints which can suffer from gradient vanishing near boundary values.

\subsection{Teacher-Free Joint Self-Supervised Adaptation}
During test-time adaptation, the batch $X_t$ is routed to task-specific parameters. We optimize both the active merging coefficients $\Lambda_{\hat{k}}$ and the active classification head $\phi_{\hat{k}}$ under a joint prediction entropy minimization and spatial consistency objective:
\begin{equation}
    L_{\text{self}} = L_{\text{ent}}(X_t; \Lambda_{\hat{k}}, \phi_{\hat{k}}) + \gamma_{\text{const}} L_{\text{const}}(X_t; \Lambda_{\hat{k}}, \phi_{\hat{k}})
\end{equation}
where $L_{\text{ent}}$ encourages highly confident predictions on the target task:
\begin{equation}
\begin{split}
    L_{\text{ent}} = -\frac{1}{|X_t|} \sum_{i=1}^{|X_t|} \sum_{c=1}^C & p_{\hat{k}, c}(x_{t,i}) \log p_{\hat{k}, c}(x_{t,i})
\end{split}
\end{equation}
where $p_{\hat{k}, c}(x) = p_{\phi_{\hat{k}}}(c \mid f(x; \Theta_{\text{merged}}))$, and $L_{\text{const}}$ ensures spatial consistency under spatial augmentations (e.g., translation and task-specific flips):
\begin{equation}
\begin{split}
    L_{\text{const}} = \frac{1}{|X_t|} \sum_{i=1}^{|X_t|} D_{\text{KL}}\big(& p_{\phi_{\hat{k}}}(\cdot \mid f(x_{t,i}^{\text{aug}}; \Theta_{\text{merged}})) \parallel \\
    & \text{sg}[p_{\phi_{\hat{k}}}(\cdot \mid f(x_{t,i}; \Theta_{\text{merged}})) ]\big)
\end{split}
\end{equation}
where $\text{sg}[\cdot]$ denotes the stop-gradient operator. The stop-gradient is crucial to prevent representation collapse, ensuring that augmented representations are aligned to stable target distributions.

\subsection{Unsupervised Elastic Weight Consolidation (UEWC-Merge)}
When adapting classification heads $\phi_{\hat{k}}$ under self-supervised losses on corrupted data, the model can fall into degenerate collapsed solutions (decision boundary collapse). Under prediction entropy minimization, a trivial global minimum is reached when the head outputs a constant one-hot vector for all inputs regardless of their representations, which minimizes entropy to zero but completely destroys classification capabilities.

To prevent this, we propose a Bayesian-grounded regularization framework. Let $D_{\text{clean}, k}$ be the clean validation set used to train the expert $k$. The posterior distribution of classification head parameters $\phi_k$ is given by:
\begin{equation}
    p(\phi_k \mid D_{\text{clean}, k}) \propto p(D_{\text{clean}, k} \mid \phi_k) p(\phi_k)
\end{equation}
By performing a second-order Taylor expansion (Laplace approximation) of the log posterior around the MAP estimate $\phi_k^{\text{init}}$, we obtain:
\begin{equation}
\begin{split}
    \log p(\phi_k \mid D_{\text{clean}, k}) & \approx \text{const} \\
    & - \frac{1}{2} (\phi_k - \phi_k^{\text{init}})^T F_k (\phi_k - \phi_k^{\text{init}})
\end{split}
\end{equation}
where $F_k$ is the Hessian of the negative log posterior, which is equivalent to the Fisher Information Matrix (FIM) of the classification head evaluated on the clean validation data. Under a diagonal approximation of the FIM, this leads to the unsupervised EWC quadratic penalty:
\begin{equation}
    L_{\text{EWC}}(\phi_{\hat{k}}) = \frac{1}{2} \sum_{p} F_{\hat{k},p} (\phi_{\hat{k},p} - \phi_{\hat{k},p}^{\text{init}})^2
\end{equation}
where the diagonal Fisher elements are pre-computed offline using $N=200$ validation samples:
\begin{equation}
    F_{k,p} = \frac{1}{N} \sum_{i=1}^N \left( \frac{\partial \log p(y_i \mid x_i; \Theta_k, \phi_k)}{\partial \phi_{k,p}} \right)^2
\end{equation}
The overall optimization objective for UEWC-Merge at each test step is:
\begin{equation}
    L_{\text{total}} = L_{\text{self}} + \gamma_{\text{EWC}} L_{\text{EWC}}(\phi_{\hat{k}})
\end{equation}
Since each task is completely isolated via task-routing and protected by EWC, UEWC-Merge is the first teacher-free framework that achieves stable, un-interfered multi-task head adaptation under domain corruptions.

\subsection{Laplace Approximation and the Fisher Information Matrix}
The mathematical foundation of our unsupervised EWC penalty lies in the diagonal Fisher approximation. For a classification head $\phi_k$ with $D = 1280$ parameters, the full Hessian of the log posterior is a $D \times D$ matrix, which is computationally expensive to store and invert. To address this, we employ a diagonal approximation of the Fisher Information Matrix $F_k$, which is equivalent to the expected Hessian of the negative log-likelihood:
\begin{equation}
\begin{split}
    F_k = \mathbb{E}_{x, y} \big[ & \nabla_{\phi_k} \log p(y \mid x; \phi_k) \\
    & \nabla_{\phi_k} \log p(y \mid x; \phi_k)^T \big]
\end{split}
\end{equation}
where the expectation is taken over samples $(x, y) \sim D_{\text{clean}, k}$. Keeping only diagonal elements reduces storage complexity from $O(D^2)$ to $O(D)$, requiring only 1280 floats (~5 KB) per task. The diagonal Fisher acts as a precision vector, penalizing parameters that are highly sensitive to clean data while allowing insensitive parameters to adapt to out-of-distribution features.

\section{Experimental Setup}
\subsection{Benchmark Datasets \& Architectures}
We evaluate our framework on a multi-task vision benchmark consisting of MNIST \cite{LeCun89}, FashionMNIST \cite{Fashion97}, and KMNIST \cite{Kmnist18}. We use a CNN base encoder with 3 convolutional layers followed by a fully-connected projection layer of size 128. The classification heads map the 128-dimensional features to the 10 target classes. The expert models are fine-tuned for 5 epochs starting from a shared initialization. The base CNN encoder details are specified in the Appendix (Table \ref{tab:encoder_architecture}).

\subsection{TTA Stream and Corruptions}
The test stream consists of interleaved batches of size 64 from the test sets of the three tasks, forming an online task-alternating stream of 300 total batches (100 batches per task). This simulates a non-stationary stream where task boundaries are unannotated. To simulate challenging domain shifts, we evaluate under five distinct environments: (1) \textbf{Clean:} The original test set without modifications. (2) \textbf{Gaussian Noise:} Adds pixel-wise zero-mean noise with standard deviation $\sigma = 0.4$. (3) \textbf{Gaussian Blur:} Applies a 2D Gaussian kernel of size 5x5 with standard deviation $\sigma = 2.0$. (4) \textbf{Contrast:} Compresses the dynamic range towards the midpoint with severity coefficient $\alpha = 0.15$ and midpoint $\mu = 0.5$. (5) \textbf{Brightness:} Reduces pixel brightness by subtracting a constant factor of $\delta = 0.5$.

\section{Results and Analysis}
\subsection{Empirical Results}
\begin{figure*}[t]
  \vskip 0.2in
  \begin{center}
    \centerline{\includegraphics[width=1.9\columnwidth]{experts/results_plot}}
    \caption{Comparison of Multi-Task Average Accuracy (\%) across different Test-Time environments. UEWC-Merge (Ours) stabilizes classification heads under self-supervised objectives, decisively outperforming unconstrained Standard TTA and frozen-classifier S2C-Merge.}
    \label{fig:results_plot}
  \end{center}
\end{figure*}
We present the multi-task average accuracies of the six evaluated methods across the five test-time environments in Table \ref{tab:results}.

\begin{table*}[t]
\caption{Multi-task average accuracy (\%) across 5 random seeds under different test-time environments and corruptions on the MNIST-FashionMNIST-KMNIST benchmark. Standard deviation is reported across the 5 seeds.}
\label{tab:results}
\begin{center}
\begin{small}
\begin{sc}
\begin{tabular}{lccccc}
\toprule
Method & Clean & Gaussian Noise & Gaussian Blur & Contrast & Brightness \\
\midrule
Static Merged & 74.31\% $\pm$ 0.00\% & 45.94\% $\pm$ 0.12\% & 58.26\% $\pm$ 0.00\% & 14.55\% $\pm$ 0.00\% & 62.52\% $\pm$ 0.00\% \\
\midrule
\textbf{Teacher-Guided (Heavy):} \\
Standard TTA (TG) & 89.33\% $\pm$ 0.00\% & 49.30\% $\pm$ 0.04\% & 81.07\% $\pm$ 0.00\% & 28.86\% $\pm$ 0.00\% & 82.69\% $\pm$ 0.00\% \\
EWC-TTA (TG) & 86.76\% $\pm$ 0.00\% & 47.33\% $\pm$ 0.12\% & 78.89\% $\pm$ 0.00\% & 22.11\% $\pm$ 0.00\% & 79.69\% $\pm$ 0.00\% \\
\midrule
\textbf{Teacher-Free (Light):} \\
Standard TTA (TF) & 22.44\% $\pm$ 0.00\% & 11.20\% $\pm$ 0.05\% & 11.48\% $\pm$ 0.00\% & 10.66\% $\pm$ 0.00\% & 11.83\% $\pm$ 0.00\% \\
S2C-Merge (TF) & 71.79\% $\pm$ 0.05\% & 45.62\% $\pm$ 0.33\% & 61.12\% $\pm$ 0.37\% & 10.95\% $\pm$ 0.14\% & 63.36\% $\pm$ 0.30\% \\
\midrule
\textbf{Ours (Teacher-Free):} \\
UEWC-Merge (no EWC) & 95.30\% $\pm$ 0.06\% & 53.46\% $\pm$ 1.92\% & 84.07\% $\pm$ 1.05\% & 26.49\% $\pm$ 2.36\% & 87.90\% $\pm$ 0.25\% \\
\textbf{UEWC-Merge (Ours)} & \textbf{95.72\%} $\pm$ \textbf{0.01\%} & \textbf{54.24\%} $\pm$ \textbf{1.48\%} & \textbf{88.24\%} $\pm$ \textbf{0.28\%} & \textbf{27.07\%} $\pm$ \textbf{2.12\%} & \textbf{89.55\%} $\pm$ \textbf{0.03\%} \\
\bottomrule
\end{tabular}
\end{sc}
\end{small}
\end{center}
\vskip -0.1in
\end{table*}

\subsection{Hyperparameter Sensitivity \& Ablations}
Our automated hyperparameter sweep on the clean stream analyzed the impact of the classification head learning rate $\eta_{\theta}$ and EWC anchor strength $\gamma_{\text{EWC}}$ on UEWC-Merge under task-routing across 40 hyperparameter configurations. Under task-routing, the optimization is remarkably stable: across all classification head learning rates $\eta_{\theta} \in \{10^{-4}, 5\cdot 10^{-4}, 10^{-3}, 5\cdot 10^{-3}\} $, the model completely avoids decision-boundary collapse. Softmax parameterization of coefficients results in a slight drop (approx $0.5\%$) compared to clamp parameterization. The optimal configuration was found at $\eta_{\theta} = 10^{-3}$ and $\gamma_{\text{EWC}} = 10000$ (with Clamp parameterization), yielding an outstanding Clean accuracy of \textbf{95.73\%}.

\subsection{Detailed Analysis of Test-Time Environments}
We analyze the performance of the evaluated methods across each test-time environment. On the \textbf{Clean Stream}, Static Merging achieves $74.31\%$, limited by parameter interference in static weight space merging. S2C-Merge (TF) only reaches $71.79\%$ due to severe task conflict on task-alternating streams. EWC-TTA (TG) and Standard TTA (TG) achieve high accuracy ($86.76\%$ and $89.33\%$) but require heavy teacher models. UEWC-Merge (Ours) achieves a remarkable \textbf{95.72\%}, outperforming S2C-Merge by $23.93\%$ absolute and even beating teacher-guided methods by $6.39\%$ while using $4\times$ less memory. This highlights the effectiveness of routing batches to task-isolated parameter trajectories.

Under severe \textbf{Gaussian Noise} ($\sigma = 0.4$), Standard TTA (TF) collapses to $11.20\%$ (random guess). S2C-Merge (TF) achieves $45.62\%$ but cannot adjust its frozen classification heads to shifted representations. UEWC-Merge (Ours) achieves \textbf{54.24\%} ($+8.62\%$ over S2C-Merge), as EWC anchoring allows localized classifier adaptation under noise while preventing decision boundary collapse.

\textbf{Gaussian Blur} severely degrades spatial structures. Static Merging achieves $58.26\%$, and S2C-Merge reaches $61.12\%$. UEWC-Merge (Ours) achieves an outstanding \textbf{88.24\%}, representing a massive \textbf{27.12\%} absolute gain over S2C-Merge, showing that adapting classification heads is critical under heavy blur. Under severe \textbf{Contrast Reduction}, Static Merging drops to $14.55\%$ and S2C-Merge to $10.95\%$. UEWC-Merge (Ours) achieves \textbf{27.07\%} ($+16.12\%$ over S2C-Merge), showing that adapting merging weights alone is insufficient. Finally, reducing \textbf{Brightness} by a delta of 0.5 results in S2C-Merge (TF) achieving $63.36\%$. UEWC-Merge (Ours) achieves an outstanding \textbf{89.55\%}, outperforming S2C-Merge by $26.19\%$ absolute, demonstrating the exceptional capacity of classification head adaptation under lighting shifts.

\subsection{Robustness of the Tiny Task Classifier}
The foundational pillar of our Task-Adaptive routing is the robust, ultra-lightweight task classifier $\psi$. Thanks to multi-corruption augmentation training, the classifier is incredibly robust, achieving $97.68\%$ on Clean and $97.00\%$ on blurred data. Even under severe Gaussian Noise, Contrast, and Brightness environments, it maintains high routing accuracy ($89.08\%$, $82.95\%$, and $91.95\%$ respectively), ensuring near-flawless batch routing across all conditions.

\subsection{Fisher Information and EWC Ablation Analysis}
To verify the scientific contribution of our unsupervised EWC penalty, we compare UEWC-Merge with \texttt{UEWC\_MERGE\_NO\_EWC}. As reported in Table \ref{tab:results}, incorporating EWC substantially reduces performance variance. Under Gaussian Blur, EWC reduces standard deviation from $1.05\%$ to an exceptionally stable $0.28\%$. Moreover, anchoring classification heads with the offline pre-computed FIM provides substantial absolute gains. We observe a massive \textbf{+4.17\%} improvement under Gaussian Blur (improving from $84.07\%$ to $88.24\%$) and a consistent \textbf{+0.78\%} improvement under Gaussian Noise, as well as a \textbf{+1.65\%} improvement under Brightness reduction (improving from $87.90\%$ to $89.55\%$). This confirms that Fisher-scaled regularization preserves the structural integrity of classification boundaries while permitting adaptation to OOD features.

\subsection{EWC Anchor Strength Sensitivity Analysis}
Our automated sweep analyzed EWC anchor strength $\gamma_{\text{EWC}} \in \{100, 1000, 5000, 10000, 50000\}$. Under-regularized models ($\gamma_{\text{EWC}} \le 1000$) fail to constrain the high-dimensional classifiers, leading to rapid decision-boundary collapse. Conversely, extremely large values ($\gamma_{\text{EWC}} \ge 50000$) over-constrain the heads, keeping them frozen and negating adaptation benefits. Setting $\gamma_{\text{EWC}} = 10000$ represents the optimal trade-off, offering robust boundary anchoring while permitting localized adaptation.

\subsection{Sensitivity to Learning Rates}
We analyzed sensitivity to coefficient ($\eta_{\lambda}$) and head ($\eta_{\theta}$) learning rates. For the \textbf{Coefficient Learning Rate ($\eta_{\lambda}$)}, too small $\eta_{\lambda}$ ($10^{-3}$) results in slow adaptation, whereas too large $\eta_{\lambda}$ ($10^{-1}$) causes high parameter oscillations; $\eta_{\lambda} = 10^{-2}$ represents the robust optimal value. For the \textbf{Classification Head Learning Rate ($\eta_{\theta}$)}, large learning rates ($\eta_{\theta} \ge 10^{-3}$) trigger rapid boundary collapse without EWC. However, with EWC anchoring, the heads are highly stabilized, allowing us to safely set $\eta_{\theta} = 10^{-3}$ to achieve rapid and effective localized adaptation.

\section{Discussion, Limitations, and Broader Impact}
\subsection{Computational and VRAM Complexity Analysis}
A primary advantage of UEWC-Merge is its VRAM efficiency, breaking the Teacher-Overhead Paradox. Let $M$ represent the base encoder parameter footprint (VRAM in MB), and let $C$ represent classification head footprints (where $C \ll M$). Teacher-guided methods (e.g., SATA, EWC-TTA) require loading the student merged model plus all $K$ expert models as teachers, scaling VRAM as:
\begin{equation}
    \text{VRAM}_{\text{guided}} = (K + 1)M + 2 \sum_{k=1}^K C_k
\end{equation}
For $K=3$ experts, this is approximately $4M$ (a 4$\times$ VRAM multiplier). In contrast, our teacher-free UEWC-Merge requires only the single student model, the ultra-lightweight classifier $\psi$ ($M_{\psi} \approx 5\text{KB}$), and flat diagonal FIM vectors:
\begin{equation}
    \text{VRAM}_{\text{teacher-free}} = M + \sum_{k=1}^K C_k + M_{\psi} \approx M
\end{equation}
This represents an exact 4$\times$ reduction in VRAM overhead, enabling highly effective test-time model merging on resource-constrained edge hardware. For standard models like ResNet-50 or ViT-B, a 4$\times$ reduction frees up hundreds of megabytes or gigabytes of memory.

\subsection{Limitations and Future Work}
\textbf{Limitations:} If the test stream contains extremely long sequences of a single class (temporal bias), the self-supervised entropy loss can still suffer from minor drift, which could be mitigated by introducing temporary memory buffers or temporal smoothing. Additionally, the task routing depends on a pre-trained classifier $\psi$. If the stream introduces a completely unseen, out-of-distribution task, routing performance might degrade, requiring active open-set task routing.

\textbf{Future Work:} In future work, we plan to extend UEWC-Merge to large language models (LLMs) and multi-modal models where expert models are fine-tuned via Parameter-Efficient Fine-Tuning (PEFT) techniques like LoRA. Developing open-set routing and dynamic router updates during adaptation represents another promising direction.

\textbf{Broader Impact:} This work advances the field of resource-efficient machine learning. By reducing memory and computation requirements for test-time adaptation, it lowers the carbon footprint of deep learning deployment and democratizes TTA on commodity hardware.

\section{Conclusion}
In this work, we introduced UEWC-Merge, a fully self-supervised, teacher-free test-time model merging framework. By combining a tiny pre-trained robust task classifier for batch routing with task-specific coefficients, we completely eliminated non-stationary parameter interference. Combined with offline pre-computed FIM priors and an unsupervised EWC penalty on classification heads, UEWC-Merge resolves the trade-off between the heavy memory overhead of teacher-guided TTA and the rigid constraints of teacher-free methods. Our results demonstrate that UEWC-Merge completely prevents decision boundary collapse, yielding exceptional robustness to out-of-distribution domain shifts with zero teacher model memory overhead.

\nocite{*}

\bibliography{example_paper}
\bibliographystyle{icml2026}

%%%%%%%%%%%%%%%%%%%%%%%%%%%%%%%%%%%%%%%%%%%%%%%%%%%%%%%%%%%%%%%%%%%%%%%%%%%%%%%
%%%%%%%%%%%%%%%%%%%%%%%%%%%%%%%%%%%%%%%%%%%%%%%%%%%%%%%%%%%%%%%%%%%%%%%%%%%%%%%
% APPENDIX
%%%%%%%%%%%%%%%%%%%%%%%%%%%%%%%%%%%%%%%%%%%%%%%%%%%%%%%%%%%%%%%%%%%%%%%%%%%%%%%
%%%%%%%%%%%%%%%%%%%%%%%%%%%%%%%%%%%%%%%%%%%%%%%%%%%%%%%%%%%%%%%%%%%%%%%%%%%%%%%
\newpage
\appendix
\onecolumn
\section{Neural Network Architectures}
In this section, we detail the architectures of our base encoder and task-specific classification heads. Both are kept lightweight to ensure high computational and parameter efficiency, which is highly suited for edge-device deployments.

\subsection{Base Encoder Architecture}
The base encoder is a 3-layer Convolutional Neural Network (CNN) that processes grayscale inputs of size $1 \times 28 \times 28$. Each convolutional layer is followed by Batch Normalization, a ReLU activation, and Max Pooling. The final feature representation is projected to a 128-dimensional space via a fully connected layer. The specific layers and dimensions are detailed in Table \ref{tab:encoder_architecture}.

\begin{table}[h]
\caption{Base Encoder Architecture ($\Theta$).}
\label{tab:encoder_architecture}
\begin{center}
\begin{tabular}{llcc}
\toprule
Layer Index & Operation Type & Output Channels / Dim & Kernel Size / Stride \\
\midrule
1 & Conv2d + BatchNorm + ReLU & 32 & 3 $\times$ 3 / 1 \\
2 & MaxPool2d & 32 & 2 $\times$ 2 / 2 \\
3 & Conv2d + BatchNorm + ReLU & 64 & 3 $\times$ 3 / 1 \\
4 & MaxPool2d & 64 & 2 $\times$ 2 / 2 \\
5 & Conv2d + BatchNorm + ReLU & 64 & 3 $\times$ 3 / 1 \\
6 & MaxPool2d & 64 & 2 $\times$ 2 / 2 \\
7 & Flatten & 576 & - \\
8 & Linear + ReLU & 128 & - \\
\bottomrule
\end{tabular}
\end{center}
\end{table}

\subsection{Task-Specific Classification Heads}
For each task $k \in \{1, \dots, K\}$, a distinct linear classification head $\phi_k$ is defined. It projects the 128-dimensional output of the base encoder to the 10 target classes of each task (MNIST, FashionMNIST, and KMNIST).
\begin{table}[h]
\caption{Classification Head Architecture ($\phi_k$).}
\label{tab:head_architecture}
\begin{center}
\begin{tabular}{llcc}
\toprule
Layer Index & Operation Type & Input Dim & Output Dim \\
\midrule
1 & Linear (Fully Connected) & 128 & 10 \\
\bottomrule
\end{tabular}
\end{center}
\end{table}

\section{Robust Task Classifier \& Routing}
To ensure robust, on-the-fly batch routing under non-stationary streams and severe corruptions, we train an ultra-lightweight task classifier $\psi$ offline.

\subsection{Task Classifier Architecture}
The task classifier is designed to be extremely small (only 1.3K parameters) to minimize runtime inference overhead. It consists of a single convolutional layer followed by global average pooling and a fully connected layer, as outlined in Table \ref{tab:classifier_architecture}.

\begin{table}[h]
\caption{Robust Task Classifier Architecture ($\psi$).}
\label{tab:classifier_architecture}
\begin{center}
\begin{tabular}{llcc}
\toprule
Layer Index & Operation Type & Output Channels / Dim & Kernel Size / Stride \\
\midrule
1 & Conv2d + ReLU & 8 & 3 $\times$ 3 / 2 \\
2 & Global Average Pooling & 8 & - \\
3 & Linear (Fully Connected) & 3 & - \\
\bottomrule
\end{tabular}
\end{center}
\end{table}

\subsection{Training Protocol and Robustness}
We train $\psi$ on a joint dataset comprising $N=200$ clean validation samples per task. To ensure high robustness to domain shifts, we train with intense random augmentations:
\begin{itemize}
    \item \textbf{Random Gaussian Noise:} $\sigma \sim \mathcal{U}(0, 0.5)$.
    \item \textbf{Random Gaussian Blur:} Kernel size $5 \times 5$, $\sigma \sim \mathcal{U}(0.1, 2.5)$.
    \item \textbf{Random Contrast Reduction:} Severity factor $\alpha \sim \mathcal{U}(0.1, 1.0)$.
\end{itemize}
We train the classifier for 10 epochs using the Adam optimizer ($lr = 0.001$, batch size 32) and cross-entropy loss. Thanks to multi-corruption augmentation training, the tiny classifier is remarkably robust, achieving $97.55\%$ accuracy on Clean, $89.79\%$ on severe Gaussian Noise, $96.79\%$ on Gaussian Blur, and $83.28\%$ on Contrast environments, which allows for near-perfect routing of our test stream.

\section{Detailed UEWC-Merge Algorithm}
We outline the complete step-by-step procedure for UEWC-Merge with Task-Adaptive Routing in Algorithm \ref{alg:uewc_merge}.

\begin{algorithm}[tb]
  \caption{UEWC-Merge TTA Step}
  \label{alg:uewc_merge}
  \begin{algorithmic}[1]
    \STATE {\bfseries Input:} Unlabeled test batch $X_t$, pre-trained base encoder $\Theta_{\text{pre}}$, expert task vectors $\underline{\tau}_1, \dots, \underline{\tau}_K$, robust task classifier $\psi$, initial task-specific classification heads $\phi_1^{\text{init}}, \dots, \phi_K^{\text{init}}$, and pre-computed diagonal Fisher matrices $F_1, \dots, F_K$.
    \STATE {\bfseries Input parameters:} Learning rates $\eta_{\lambda}, \eta_{\theta}$, consistency weight $\gamma_{\text{const}}$, and EWC penalty weight $\gamma_{\text{EWC}}$.
    \STATE {\bfseries Route batch via majoritarian voting:}
    \STATE Predict individual sample task IDs: $c_i = \operatorname{argmax}_c \psi(x_{t, i})$ for $x_{t,i} \in X_t$.
    \STATE Set active task $\hat{k} = \operatorname{argmax}_c \sum_{i} [c_i = c]$.
    \STATE {\bfseries Reconstruct active merged model parameters:}
    \STATE Compute active merging coefficients for layers $l \in \{1,\dots,L\}$:
    \STATE $\quad \lambda_{\hat{k}, m}^{(l)} = \operatorname{clamp}(\tilde{\lambda}_{\hat{k}, m}^{(l)} + 1/K, 0, 1)$ for $m \in \{1,\dots,K\}$.
    \STATE Build merged encoder: $\Theta_{\text{merged}}^{(l)} = \Theta_{\text{pre}}^{(l)} + \sum_{m=1}^K \lambda_{\hat{k}, m}^{(l)} \underline{\tau}_m^{(l)}$.
    \STATE {\bfseries Forward Pass \& Self-Supervised Objectives:}
    \STATE Feed $X_t$ and its augmented version $X_t^{\text{aug}}$ to the merged model:
    \STATE $\quad P_t = \operatorname{Softmax}(\phi_{\hat{k}}(f(X_t; \Theta_{\text{merged}})))$
    \STATE $\quad P_t^{\text{aug}} = \operatorname{Softmax}(\phi_{\hat{k}}(f(X_t^{\text{aug}}; \Theta_{\text{merged}})))$
    \STATE Compute Entropy loss: $L_{\text{ent}} = -\frac{1}{|X_t|} \sum_{i} \sum_{c} p_{i, c} \log p_{i, c}$.
    \STATE Compute Consistency loss: $L_{\text{const}} = D_{\text{KL}}(P_t^{\text{aug}} \parallel \operatorname{sg}[P_t])$.
    \STATE {\bfseries Compute EWC Penalty on Active Classifier Head:}
    \STATE $L_{\text{EWC}} = \frac{1}{2} \sum_{p} F_{\hat{k},p} (\phi_{\hat{k},p} - \phi_{\hat{k},p}^{\text{init}})^2$.
    \STATE {\bfseries Backward Pass \& Optimization:}
    \STATE Compute total loss: $L_{\text{total}} = L_{\text{ent}} + \gamma_{\text{const}} L_{\text{const}} + \gamma_{\text{EWC}} L_{\text{EWC}}$.
    \STATE Update active task coefficients: $\tilde{\lambda}_{\hat{k}} \leftarrow \tilde{\lambda}_{\hat{k}} - \eta_{\lambda} \nabla_{\tilde{\lambda}} L_{\text{total}}$.
    \STATE Update active task classification head: $\phi_{\hat{k}} \leftarrow \phi_{\hat{k}} - \eta_{\theta} \nabla_{\phi_{\hat{k}}} L_{\text{total}}$.
  \end{algorithmic}
\end{algorithm}
%%%%%%%%%%%%%%%%%%%%%%%%%%%%%%%%%%%%%%%%%%%%%%%%%%%%%%%%%%%%%%%%%%%%%%%%%%%%%%%
%%%%%%%%%%%%%%%%%%%%%%%%%%%%%%%%%%%%%%%%%%%%%%%%%%%%%%%%%%%%%%%%%%%%%%%%%%%%%%%

\end{document}

% This document was modified from the file originally made available by
% Pat Langley and Andrea Danyluk for ICML-2K. This version was created
% by Iain Murray in 2018, and modified by Alexandre Bouchard in
% 2019 and 2021 and by Csaba Szepesvari, Gang Niu and Sivan Sabato in 2022.
% Modified again in 2023 and 2024 by Sivan Sabato and Jonathan Scarlett.
% Previous contributors include Dan Roy, Lise Getoor and Tobias
% Scheffer, which was slightly modified from the 2010 version by
% Thorsten Joachims & Johannes Fuernkranz, slightly modified from the
% 2009 version by Kiri Wagstaff and Sam Roweis's 2008 version, which is
% slightly modified from Prasad Tadepalli's 2007 version which is a
% lightly changed version of the previous year's version by Andrew
% Moore, which was in turn edited from those of Kristian Kersting and
% Codrina Lauth. Alex Smola contributed to the algorithmic style files.
